\documentclass[10pt,a4paper]{article}

\usepackage[margin=1.35cm]{geometry}
\usepackage[T1]{fontenc}
\usepackage[utf8]{inputenc}
\usepackage{lmodern}
\usepackage{microtype}
\usepackage{xcolor}
\usepackage{enumitem}
\usepackage{titlesec}
\usepackage[hidelinks,unicode]{hyperref}
\usepackage{ragged2e}

\definecolor{accent}{HTML}{17365D}
\hypersetup{
  pdftitle={Talal Alrawajfeh - Senior Applied Scientist},
  pdfauthor={Talal Alrawajfeh},
  pdfsubject={Resume},
  pdfkeywords={machine learning, artificial intelligence, applied scientist, data scientist, AI engineer}
}

\pagestyle{empty}
\setlength{\parindent}{0pt}
\setlength{\parskip}{0pt}
\setlist[itemize]{leftmargin=1.25em,itemsep=1.4pt,topsep=2.5pt,parsep=0pt,partopsep=0pt}
\titleformat{\section}{\large\bfseries\color{accent}}{}{0pt}{}[\vspace{-2pt}\color{accent}\titlerule]
\titlespacing*{\section}{0pt}{7pt}{4pt}
\urlstyle{same}

\newcommand{\role}[4]{%
  \textbf{#1} \hfill \textbf{#2}\\[-1pt]
  \textit{#3} \hfill \textit{#4}\par\vspace{-2pt}
}

\begin{document}
\RaggedRight

\begin{center}
  {\LARGE\bfseries TALAL ALRAWAJFEH}\\[2pt]
  {\normalsize\bfseries Senior Applied Scientist}\\[4pt]
  Helsinki, Finland \textbar{} +358 45 253 0012 \textbar{}
  \href{mailto:talalalrawajfeh@gmail.com}{talalalrawajfeh@gmail.com}\\[2pt]
  \href{https://www.linkedin.com/in/talal-alrawajfeh}{linkedin.com/in/talal-alrawajfeh} \textbar{}
  \href{https://github.com/talalalrawajfeh}{github.com/talalalrawajfeh} \textbar{}
  \href{https://orcid.org/0000-0002-5190-4339}{ORCID: 0000-0002-5190-4339}
\end{center}
\vspace{-5pt}

\section{Professional Summary}
Researcher and applied scientist with 10 years of experience across AI research, production machine learning, and software engineering.
Five years of experience turning ambiguous, business-critical problems into production AI systems for regulated banking.
Delivered a document signature verification model with 99\% accuracy, cheque OCR/ICR automating 80\% of manual data entry, and eKYC components supporting the automation of 70\% of online onboarding applications. 
Combines mathematical depth with hands-on engineering and clear cross-functional technical leadership.

\section{Core Skills}
\textbf{Areas of Expertise:} Deep learning, machine learning, computer vision, time-series analysis and forecasting, anomaly detection, probabilistic modelling, Bayesian statistics, continual learning, neural processes, tabular foundation models, differential privacy, interpretability and explainability, Mathematics, formal verification, software engineering \\
\textbf{Programming and Tools:} Python, C, C++, x86-64 Assembly, Java, Kotlin, Scala, C\#, Ruby, SQL, Lean, PyTorch, JAX, TensorFlow, Keras, scikit-learn, NumPy, SciPy, pandas, OpenCV, Pyro, Optuna, Numba\\
\textbf{DevOps and MLOps:} MLflow, ONNX, OpenVINO, TensorFlow Lite, Docker, Kubernetes, Git, CI/CD \\
\textbf{Software Engineering:} Spring Boot, .NET Core, Cython, FastAPI, Flask, Vert.x, Rails, Node.js, Kafka, CMake, Maven, Gradle, MySQL, Oracle Database, Redis, Elasticsearch

\section{Professional Experience}
\role{Doctoral Researcher}{Feb 2024--Present}{University of Helsinki}{Helsinki, Finland}
\begin{itemize}
  \item Lead theory-to-experiment research in privacy-preserving machine learning spanning differential privacy, Bayesian inference, variational inference, continual learning, neural processes, and tabular foundation models.
  \item Introduced PrivTab, a tabular foundation model for differentially private classification that adapts to new datasets through in-context learning in a single forward pass; first-author preprint, 2026.
  \item Introduced noise-aware differentially private variational inference for high-dimensional, non-conjugate models; designed a more accurate posterior evaluation method and demonstrated calibrated results on Bayesian regression tasks. Published first-author work at AISTATS 2025.
  \item Co-led an equal-contribution TMLR 2026 paper identifying classifier output label spaces as a privacy side channel; formalized differential privacy for continual learning and evaluated two mechanisms that restore valid privacy guarantees.
  \item Formalized differential privacy mechanisms and sensitivity proofs in Lean 4; run reproducible experiments at scale using Python, JAX, PyTorch, and high-performance computing infrastructure.
  \item Teaching assistant for Introduction to Data Science (three academic years) and Practical Machine Learning.
\end{itemize}

\role{Senior Machine Learning Engineer}{May 2019--Feb 2024}{ProgressSoft Corporation}{Amman, Jordan}
\begin{itemize}
  \item Led multiple data science teams within the AI department, coordinating delivery through regular technical reviews and checkpoints while owning problem framing, experimentation, evaluation, production integration, and stakeholder communication.
  \item Led deep-learning development for document signature segmentation and offline signature verification with 99\% accuracy; delivered OCR/ICR models that automated at least 80\% of manual data entry.
  \item Built document segmentation, passive livenes, and facial recognition components for bank eKYC software, including ID-card/passport detection, perspective correction, template matching, and field extraction; contributed to automating at least 70\% of onboarding applications.
  \item Created domain-agnostic cash-flow forecasting systems combining statistical, ML, and deep-learning models; designed adaptive transaction anomaly detection that tracks customer-behaviour drift with controllable sensitivity.
  \item Deployed models as part of ProgressSoft software running on premises in banks; optimized inference with ONNX, OpenVINO, TensorFlow Lite, hardware benchmarking, and modern C++ implementations using SIMD/AVX.
  \item Technically led MessageDepot, an LLM-enabled system developed by team members for natural-language-to-SQL querying, result summarization, and chart generation.
  \item Interviewed candidates, screened CVs, onboarded employees, and created supervised self-learning improvement plans with structured checkpoints, progress reviews, and task feedback; helped coordinate work across the AI department.
\end{itemize}

\role{Backend Engineer}{Apr 2018--Apr 2019}{Strategic Directory for Shopping Sites / Blue Kangaroo}{Amman, Jordan}
\begin{itemize}
  \item Built data collection and analysis pipelines and production APIs for a white-label, AI-enabled e-commerce shopping assistant using Ruby on Rails, Kafka, and Elasticsearch.
  \item Improved reliability and maintainability through refactoring, code review, unit and integration testing, and documentation of APIs and integrations while collaborating across multiple teams.
\end{itemize}

\role{Associate Java Developer}{Sep 2016--Mar 2018}{ProgressSoft Corporation}{Amman, Jordan}
\begin{itemize}
  \item Developed and maintained large-scale electronic payment solutions for banks and corporates, including an electronic cheque-deposit service with multiple external REST and RPC integrations.
  \item Delivered regulated financial software using Java, Spring Boot, Hibernate/JPA, TDD, clean code, Agile, DevOps, performance benchmarking.
\end{itemize}

\section{Selected Publications and Research}
\begin{itemize}
  \item \textbf{Interpretable Differentially Private Continual Learning with Prototype-Based Classifiers.} Yuan He, \textbf{Talal Alrawajfeh}, Wachirapatra Mahatubgris, Agil Aghasanli, Mikko A. Heikkilä, Plamen P Angelov, Antti Honkela. Submitted to ICLR.
  \item \textbf{Efficient Provably Private Classification with a Tabular Foundation Model.} \textbf{Talal Alrawajfeh}, Cristiana Diaconu, Ossi Räisä, Sebastian Rodriguez Beltran, Yuan He, John Bronskill, Richard E. Turner, Antti Honkela. arXiv preprint, 2026. \href{https://arxiv.org/abs/2610.10068}{arXiv:2610.10068} \textbar{} \href{https://arxiv.org/pdf/2610.10068}{Paper} \textbar{} \href{https://github.com/TrustworthyMLHelsinki/PrivTab}{Code}
  \item \textbf{Privacy Leakage via Output Label Space and Differentially Private Continual Learning.} Marlon Tobaben*, \textbf{Talal Alrawajfeh}*, Marcus Klasson, Mikko A. Heikkilä, Arno Solin, Antti Honkela. TMLR, 2026. *Equal contribution. \href{https://openreview.net/forum?id=ejnVWfknCm}{Paper} \textbar{} \href{https://github.com/TrustworthyMLHelsinki/dp-continual-learning}{Code}
  \item \textbf{Noise-Aware Differentially Private Variational Inference.} \textbf{Talal Alrawajfeh}, Joonas Jälkö, Antti Honkela. AISTATS 2025, PMLR 258:4987--4995; poster presentation. \href{https://proceedings.mlr.press/v258/alrawajfeh25a.html}{Paper} \textbar{} \href{https://github.com/DPBayes/noise-aware-dpvi}{Code}
\end{itemize}

\section{Education}
\role{PhD in Science}{Feb 2024--Expected 2027}{University of Helsinki}{Helsinki, Finland}
Specialization: Computer Science

Research focus: trustworthy machine learning, privacy-preserving machine learning, Bayesian statistics, amortized inference, continual learning, neural processes, and tabular foundation models. 

Coursework average: \textbf{5.0/5.0}.

\vspace{3pt}
\role{MSc in Mathematics}{Sep 2017--Aug 2019}{University of Jordan}{Amman, Jordan}
GPA: \textbf{4.0/4.0; first in class}. 

Thesis: \textit{On D-spaces and Covering Properties}; research in general topology, covering properties, and category-theoretic approaches to the D-space problem.

\vspace{3pt}
\role{BSc in Mathematics}{Sep 2013--Aug 2016}{University of Jordan}{Amman, Jordan}
GPA: \textbf{3.83/4.0; first in class}. Completed the degree in three years.
\section{Additional Information}
\textbf{Work status:} Authorized to work in Finland; available with two weeks' notice; seeking Helsinki-based hybrid or remote roles\\
\textbf{Languages:} Arabic (native); English (full professional proficiency; TOEFL iBT 105/120); Finnish (approximately CEFR A2; University of Helsinki coursework through \href{https://studies.helsinki.fi/kurssit/toteutus/otm-67570c2a-f9ed-467f-92a7-d78330209a79}{Finnish 5A}; actively studying)\\

\end{document}
